\section{Introduction}
\IEEEPARstart{A}{gentic} AI changes the unit of clinical evaluation. A medical language model may produce a diagnosis; an agent may decide which chart to inspect, request additional evidence, invoke an imaging tool, draft an order, and attempt to change a destination system. The correctness of the final sentence cannot establish the correctness of this entire sequence. Medical language-model and foundation-model research demonstrates substantial informational capability while also identifying limitations in clinical reasoning and human interaction \cite{singhal2023,moor2023,hager2024,goh2024}. The question for healthcare agents is what additional evidence justifies their ability to cause effects.

Agency predates language models. Classical accounts emphasize situated, goal-directed interaction \cite{wooldridge1995}. Recent architectures combine reasoning and acting, learned tool use, retrieval, and feedback memory \cite{react2023,toolformer2023,rag2020,reflexion2023}. These components can make otherwise inaccessible information available and support complex workflows. They also introduce new failure surfaces: selecting the wrong patient, reading obsolete evidence, reusing an approval after its context changes, and reporting success before a destination confirms completion. A competent planner may still issue an unauthorized action. A correctly authorized operation may still be clinically inappropriate.

Several reviews already characterize healthcare agents. A scoping review examines applications and conceptual distinctions \cite{collaco2026review}; broader surveys organize agents, architectures, evaluation, and translation challenges \cite{xu2026survey,vatsal2026taxonomy,hu2026survey}. This literature precludes a credible claim that healthcare agent review or governance is unexplored. Our contribution is a complementary analysis organized around \emph{action assurance}: a proposed operation, an observed effect, and the institutional authority that permits that effect. Table~\ref{tab:position} makes the intended distinction explicit without assuming that prior reviews omit safety or tool use.

We ask three questions. First, what kind of agency and destination effect do empirical reports actually evaluate? Second, which failure mechanisms emerge between a plausible proposal and a verified, permissible effect? Third, how should software controls and clinical evaluation be connected before agents receive bounded operational authority? We address these questions through a purposive evidence map and propose Clinical Action Assurance (CAA), a reference framework with a synthetic implementation. CAA is a synthesis and operationalization of existing principles, not a certification scheme or a proof that an agent is safe. The Outstanding Questions box identifies priorities for clinical translation.

\begin{table}[t]
\caption{Positioning relative to existing review literature}
\label{tab:position}\centering\footnotesize
\begin{tabularx}{\columnwidth}{@{}p{.23\columnwidth}Y@{}}\toprule
Source & Complementary emphasis \\ \midrule
Collaco et al. \cite{collaco2026review} & Scoping of applications and agentic conceptual features.\\
Xu et al. \cite{xu2026survey} & Broad survey of AI agents in healthcare.\\
Vatsal et al. \cite{vatsal2026taxonomy} & Taxonomic treatment of healthcare agent capabilities.\\
Abdollahi et al. \cite{hu2026survey} & Biomedical engineering survey of agents and translation.\\
Present review & Action/effect/authority distinctions; traceable evidence limits; approval and destination-state obligations.\\ \bottomrule
\end{tabularx}
\end{table}
